\section{Introduction}
The useful outcome of memory expansion is additional work completed
within a time budget. A larger memory pool may accommodate more KV state
while its transaction machinery cannot sustain the associated access
rate. Conversely, data that occupies substantial space but is rarely
accessed may be a favorable candidate for remote placement. We ask:
\emph{when does Ethernet memory pooling increase the inference working set
that a fixed accelerator deployment can actually serve?}

We use MPU to denote an independently attached memory-service node.
The proposed MPU provides remote storage and access service; it does not
execute attention or expert computation. Our scope is a multi-die
accelerator system where memory traffic and inference communication may
share the Ethernet scale-up fabric.

Several access patterns make nominal capacity an incomplete metric.
Active KV reads depend on the attention algorithm and selected state.
Paused sessions or cached prefixes can remain idle and later trigger
bursty restoration. Fine-grained tables can require many transactions
despite modest byte volume. These patterns must be distinguished:
older KV entries in dense attention are not automatically cold, and
restoring a paused session is not equivalent to reading its state on
every decode step.

\subsection{Research Hypothesis}
We hypothesize that capacity-only or nominal-bandwidth placement can
overestimate admissible load when memory-service slots, endpoint windows,
shared paths, or response consumption become limiting resources. A
placement and admission policy that accounts for these constraints may
increase the workload region satisfying an explicit time budget. The
hypothesis can fail: if calibrated workloads are primarily bandwidth
limited, a simple bandwidth-proportional policy may suffice.

The intended contributions are (1) a characterization of serviceable
capacity under specified access patterns, (2) an implementable placement
and transaction-admission mechanism, and (3) a calibrated evaluation of
its benefits and failure regions. None is claimed as completed in this
draft. In particular, adding more MPU nodes is not by itself evidence
for the mechanism: it also adds ports, service bandwidth, and slots.

\subsection{Positioning Against Existing Work}
Mooncake already combines a disaggregated KV cache with SLO-oriented
scheduling and prediction-based early rejection~\cite{mooncake}.
MemServe combines elastic memory pooling, context caching, and
locality-aware scheduling~\cite{memserve}. Fang and colleagues formulate
dynamic KV placement under capacity constraints to exploit heterogeneous
memory bandwidth~\cite{dynamicplacement}. Consequently, pooling,
capacity-aware placement, and admission control are not sufficient
novelty claims. The proposed distinction is the translation of concrete
Ethernet MPU transaction and response-path constraints into placement
and admission decisions. A full related-work review and experimental
comparison are still required to establish that distinction.
